% =====================================================
% 题型：钝角三角形面积与边长选择题 (q_triangle_obtuse_area_side.tex)
% =====================================================
% 难度：★★ (中档)
% =====================================================
\newcommand{\qTriangleObtuseAreaSideStem}{在 \(\triangle ABC\) 中，角 \(A,B,C\) 的对边分别为 \(a,b,c\)。若 \(a=3,\ b=5\)，\(C\) 为钝角，且 \(\triangle ABC\) 的面积为 \(\dfrac{15\sqrt3}{4}\)，则 \(c=\hfill（\quad）\)}

\newcommand{\qTriangleObtuseAreaSideOptions}{%
  \mcoptions[4]{\(4\)}{\(\sqrt{19}\)}{\(7\)}{\(\sqrt{61}\)}%
}

\newcommand{\qTriangleObtuseAreaSideAnswer}{D}

\newcommand{\qTriangleObtuseAreaSideSolution}{%
  由面积公式：
  \[
    \frac12ab\sin C=\frac{15\sqrt3}{4}.
  \]
  代入 \(a=3,\ b=5\)，得
  \[
    \frac{15}{2}\sin C=\frac{15\sqrt3}{4}
    \implies \sin C=\frac{\sqrt3}{2}.
  \]
  因为 \(C\) 为钝角，所以
  \[
    C=\frac{2\pi}{3},
    \qquad
    \cos C=-\frac12.
  \]
  由余弦定理：
  \[
    c^2=a^2+b^2-2ab\cos C
    =9+25-30\left(-\frac12\right)
    =49.
  \]
  因此 \(c=7\)。\\
  故选 C。%
}

\newcommand{\qTriangleObtuseAreaSideDiagnosis}{%
  高频错误是由 \(\sin C=\frac{\sqrt3}{2}\) 直接写成 \(C=\frac\pi3\)，忽略题目已明确 \(C\) 为钝角。%
}

\newcommand{\qTriangleObtuseAreaSideTeachingNotes}{%
  题目链条非常标准：面积公式求 \(\sin C\) → 结合钝角条件确定角 → 余弦定理求边。重点训练“一个正弦值对应两个互补角”的分支意识。%
}
